% Author: Prof. Dr. Matthias Jung, DL9MJ
% Year: 2026
%
% Vertikaldiagramme idealer Monopolantennen über perfekter Erde
%
% Elevationswinkel alpha:
%   0°  = Horizont
%   90° = Zenit
%
% Gezeigt werden:
%   lambda/4
%   lambda/2
%   5 lambda/8
%   lambda
%
% Die orangen Kurven zeigen das normierte Fernfeld.
% Die roten Kurven zeigen die idealisierte Stromverteilung
% entlang des Strahlers.
% Referenz
% https://en.wikipedia.org/wiki/Monopole_antenna#mwBYI

\begin{tikzpicture}

% --------------------------------------------------------------------------
% Fernfeld eines idealen vertikalen Monopols
%
% alpha = Elevationswinkel über dem Horizont
% h     = Strahlerlänge in Wellenlängen
%
% E(alpha) ~
%
% | cos(2*pi*h*sin(alpha)) - cos(2*pi*h) |
% -------------------------------------------------
%                  cos(alpha)
%
% PGF rechnet trigonometrische Funktionen in Grad.
% Daher entspricht 2*pi*h dem Winkel 360*h.
% --------------------------------------------------------------------------

\pgfmathdeclarefunction{monopat}{3}{%
    \pgfmathparse{%
        abs(
            (cos(360*#2*sin(#1))-cos(360*#2))
            / cos(#1)
        ) / #3
    }%
}

% Radius eines Polardiagramms
\def\R{1.40}

% Horizontale Auslenkung der Stromkurve
\def\Iscale{0.28}

% Vertikale Skalierung der Antennendarstellung
\def\AntennaScale{1.55}

\tikzset{
    polargrid/.style={
        draw=black!35,
        line width=0.35pt
    },
    polargridminor/.style={
        draw=black!20,
        line width=0.25pt
    },
    radiation/.style={
        draw=DARCorange,
        line width=1.4pt
    },
    antenna/.style={
        draw=black,
        line width=1.1pt
    },
    current/.style={
        draw=DARCred,
        line width=1.3pt
    }
}

% --------------------------------------------------------------------------
% Panel
%
% #1 x-Position
% #2 y-Position
% #3 Strahlerlänge h/lambda
% #4 Normierung des Feldes
% #5 Beschriftung der mechanischen Länge
% --------------------------------------------------------------------------

\newcommand{\MonopolePanel}[5]{%
    \begin{scope}[shift={(#1,#2)}]

        % --------------------------------------------------------------
        % Polardiagramm
        % --------------------------------------------------------------

        % konzentrische Bögen
        \foreach \r in {0.2,0.4,0.6,0.8,1.0}{
            \draw[polargridminor]
                (0,0)
                ++(0:{\r*\R})
                arc[
                    start angle=0,
                    end angle=90,
                    radius={\r*\R}
                ];
        }

        % Winkelstrahlen
        \foreach \a in {0,10,...,90}{
            \draw[polargrid]
                (0,0) -- (\a:\R);
        }

        % Winkelbeschriftungen
        \foreach \a in {0,10,...,90}{
            \node[
                font=\scriptsize,
                inner sep=1pt
            ] at (\a:{1.10*\R}) {\a};
        }

        % --------------------------------------------------------------
        % Strahlungsdiagramm
        % --------------------------------------------------------------

        \draw[radiation]
            plot[
                domain=0:89.7,
                samples=350,
                smooth,
                variable=\a
            ]
            (
                {\R*monopat(\a,#3,#4)*cos(\a)},
                {\R*monopat(\a,#3,#4)*sin(\a)}
            );

        % --------------------------------------------------------------
        % Antennendarstellung
        % --------------------------------------------------------------

        \pgfmathsetmacro{\antennaheight}{\AntennaScale*#3}

        % Groundplane / Erde
        \draw[antenna]
            (-0.73,0) -- (-0.27,0);

        % Strahler
        \draw[antenna]
            (-0.50,0) -- (-0.50,\antennaheight);

        % --------------------------------------------------------------
        % Stromverteilung
        %
        % I(z) = Imax * sin(k * (l-z))
        %
        % Da z und l in Wellenlängen angegeben sind:
        %
        % I(z) = sin(360° * (l-z))
        %
        % Positive und negative Werte werden auf unterschiedlichen
        % Seiten des Strahlers dargestellt und zeigen damit auch
        % Phasenumkehrungen des Stroms.
        % --------------------------------------------------------------

        \draw[current]
            plot[
                domain=0:#3,
                samples=180,
                smooth,
                variable=\z
            ]
            (
                {
                    -0.50
                    + \Iscale*sin(360*(#3-\z))
                },
                {
                    \AntennaScale*\z
                }
            );

        % --------------------------------------------------------------
        % Längenmaß
        % --------------------------------------------------------------

        \draw[
            Triangle-Triangle,
            line width=0.55pt
        ]
            (-0.82,0)
            --
            (-0.82,\antennaheight);

        % Längenbeschriftung
        \node[
            rotate=90,
            font=\small,
            anchor=south
        ] at (-0.87,{0.5*\antennaheight})
            {$l=#5$};

    \end{scope}
}

% --------------------------------------------------------------------------
% Obere Reihe
% --------------------------------------------------------------------------

% lambda/4
%
% Strom:
% Maximum am Speisepunkt
% Nullstelle an der Spitze
%
% Strahlung:
% Maximum am Horizont
%
\MonopolePanel
    {0}
    {0}
    {0.25}
    {1.0}
    {\frac{1}{4}\lambda}


% lambda/2
%
% Strom:
% Nullstelle am Speisepunkt
% Maximum bei lambda/4
% Nullstelle an der Spitze
%
% Strahlung:
% Maximum am Horizont
%
\MonopolePanel
    {3.15}
    {0}
    {0.50}
    {2.0}
    {\frac{1}{2}\lambda}


% --------------------------------------------------------------------------
% Untere Reihe
% --------------------------------------------------------------------------

% 5 lambda/8
%
% Strom:
% kleiner Stromabschnitt mit umgekehrter Phase am Speisepunkt
% Stromnullstelle bei z = lambda/8
% großer Strombauch darüber
%
% Strahlung:
% Hauptmaximum am Horizont
% erste Nullstelle bei ca. 36.87°
% kleine Hochwinkelkeule bei ca. 58.91°
%
\MonopolePanel
    {0}
    {-2.05}
    {0.625}
    {1.707106781}
    {\frac{5}{8}\lambda}


% lambda
%
% Strom:
% Nullstellen bei
% z = 0
% z = lambda/2
% z = lambda
%
% zwei Strombäuche mit entgegengesetzter Phase
%
% Strahlung:
% Nullstelle am Horizont
% Hauptmaximum bei ca. 32.56°
%
\MonopolePanel
    {3.15}
    {-2.05}
    {1.0}
    {2.338981541}
    {\lambda}

\end{tikzpicture}